\documentclass[11pt,a4paper]{article}
\usepackage[T1]{fontenc}
\usepackage[utf8]{inputenc}
\usepackage{lmodern}
\usepackage[margin=27mm,headheight=15pt]{geometry}
\usepackage{amsmath,amssymb,amsthm,mathtools}
\usepackage{microtype,booktabs,longtable,array,enumitem}
\usepackage{xcolor,fancyhdr}
\usepackage[hidelinks,breaklinks]{hyperref}
\usepackage{listings}
\definecolor{ink}{HTML}{19344D}
\definecolor{muted}{HTML}{56626B}
\hypersetup{pdftitle={Brun's theorem in intuitionistic primitive recursive arithmetic},
 pdfauthor={Research prepared for Vladimir Reshetnikov},
 pdfsubject={Finite Selberg sieve, primitive-recursive Cauchy modulus, Lean roadmap}}
\pagestyle{fancy}
\fancyhf{}
\fancyhead[L]{\small\color{muted}Brun's theorem in iPRA}
\fancyhead[R]{\small\color{muted}Finite proof and formalization plan}
\fancyfoot[C]{\thepage}
\renewcommand{\headrulewidth}{0.3pt}
\setlength{\parskip}{4pt}
\setlength{\parindent}{1.2em}
\setlist{itemsep=3pt,topsep=5pt}
\emergencystretch=2em
\newtheorem{theorem}{Theorem}[section]
\newtheorem{lemma}[theorem]{Lemma}
\newtheorem{proposition}[theorem]{Proposition}
\newtheorem{corollary}[theorem]{Corollary}
\theoremstyle{definition}
\newtheorem{definition}[theorem]{Definition}
\newtheorem{remark}[theorem]{Remark}
\newcommand{\N}{\mathbb N}
\newcommand{\Z}{\mathbb Z}
\newcommand{\Q}{\mathbb Q}
\newcommand{\ipra}{\mathsf{iPRA}}
\newcommand{\Twin}{\operatorname{Twin}}
\newcommand{\rad}{\operatorname{rad}}
\newcommand{\lcm}{\operatorname{lcm}}
\newcommand{\sfD}{\mathcal D}
\newcommand{\PR}{\mathsf{PR}}
\newcommand{\eps}{\varepsilon}
\newcommand{\code}[1]{\texttt{\small #1}}
\lstset{basicstyle=\ttfamily\small,breaklines=true,columns=fullflexible,
 keepspaces=true,frame=single,rulecolor=\color{muted},
 backgroundcolor=\color{black!2},showstringspaces=false}

\title{\color{ink}\bfseries Brun's theorem in intuitionistic\\
primitive recursive arithmetic\\[5pt]
\large A finite sieve proof with an explicit Cauchy modulus}
\author{Research prepared for Vladimir Reshetnikov}
\date{15 November 2026}

\begin{document}
\maketitle
\begin{abstract}
We give a fully developed finite-arithmetic proof of the convergence of the
twin-prime reciprocal series in intuitionistic primitive recursive arithmetic.
If \(T(X)\) counts primes \(p\le X\) for which \(p+2\) is prime, we prove
\[
 k^2T(2^k)\le 1400\,2^k\le 2048\,2^k.
\]
Consequently every finite reciprocal tail with smaller prime at least \(2^m\)
is at most \(8192/m\). For the rational partial sums \(B(N)\), indexed by
the smaller prime \(p<N\), the primitive recursive function
\[
 N(q)=2^{\,2^{q+13}}
\]
is a Cauchy modulus: \(M,L\ge N(q)\) implies
\(\lvert B(M)-B(L)\rvert\le2^{-q}\).

The proof constructs rational Selberg weights and proves their quadratic
identity directly. Its denominator estimate follows by grouping a finite
square of a harmonic sum according to prime support. All summations,
factorizations and searches are bounded; the induction invariants admit
primitive recursive Boolean encodings. A denominator-cleared first-order
statement makes the precise iPRA conclusion explicit. No analytic limit
theorem is used to obtain convergence.

The archive includes exact-arithmetic diagnostics, a Lean specification
with conditional assembly proofs, and a detailed route to a complete Lean
derivation in the existing ConciseIPRA object theory. The human proof is
complete; the full Lean sieve proof and object-theory derivation remain
to be implemented. The supplied Lean source has not been compiled.
\end{abstract}
\clearpage
\begingroup
\small
\setlength{\parskip}{0pt}
\tableofcontents
\endgroup
\newpage

\section{The theorem and its constructive meaning}
\label{sec:statement}

Brun's theorem asserts that the reciprocals in all twin-prime pairs have
a finite sum. We index a pair by its smaller member. Write
\[
 \Twin(p)\quad\Longleftrightarrow\quad
 p\text{ and }p+2\text{ are prime},
\]
and define, for natural numbers \(X,N\),
\[
 T(X)=\#\{p:1\le p\le X,\ \Twin(p)\},\qquad
 B(N)=\sum_{\substack{1\le p<N\\\Twin(p)}}
       \left(\frac1p+\frac1{p+2}\right).
\]
The convention deliberately counts \(5\) twice: once in \((3,5)\) and once
in \((5,7)\). Thus \(B(4)=8/15\) and \(B(6)=92/105\).
The endpoint in \(T\) is inclusive; the endpoint in \(B\) is strict.
The condition \(p+2\le X\) is \emph{not} imposed in the definition of \(T(X)\).

\begin{theorem}[Finite quantitative Brun theorem]
\label{thm:main}
The following assertions have proofs in intuitionistic primitive recursive
arithmetic, when rational notation is read through positive-denominator
integer codes.
\begin{enumerate}[label=(\roman*)]
\item For every \(k\ge0\),
\begin{equation}
 k^2T(2^k)\le 1400\,2^k\le 2^{11}2^k.
 \label{eq:main-count}
\end{equation}
\item For all \(m\ge1\) and \(L\ge0\),
\begin{equation}
 \sum_{\substack{2^m\le p<L\\\Twin(p)}}
 \left(\frac1p+\frac1{p+2}\right)
 \le \frac{2^{13}}m.
 \label{eq:main-tail}
\end{equation}
\item Define \(N(q)=2^{\,2^{q+13}}\). For all \(q,M,L\in\N\),
\begin{equation}
 M,L\ge N(q)\quad\Longrightarrow\quad
 \lvert B(M)-B(L)\rvert\le 2^{-q}.
 \label{eq:main-cauchy}
\end{equation}
\end{enumerate}
\end{theorem}

The constants are intentionally generous. They simplify the arithmetic
proof and give a completely specified modulus; they are not intended for
efficient numerical evaluation of Brun's constant. In particular, even the
cutoff at \(q=0\) is \(2^{8192}\).

\subsection{A statement literally in the language of iPRA}

There are no real-number variables in first-order iPRA. Accordingly, the
primary conclusion is a theorem about finite integer computations. Here is
one exact version.

Let \(\chi(p)\in\{0,1\}\) be the primitive recursive characteristic function
of \(\Twin(p)\). Primality requires \(p\ge2\) and is tested by trial division
over the explicitly bounded range \(2,\ldots,p-1\). Put
\[
 D(V)=(V+1)!,
\]
and define the integer
\begin{equation}
 R(U,V)=
 \sum_{U\le p<V}\chi(p)
 \left(D(V)\mathbin{\mathrm{div}}p+
       D(V)\mathbin{\mathrm{div}}(p+2)\right).
 \label{eq:integer-tail}
\end{equation}
Use total division, with a fixed value such as \(0\) at divisor \(0\).
The terms at \(p=0,1\) have coefficient \(0\). Every active denominator is
positive and at most \(V+1\), and hence divides \(D(V)\).
It follows exactly, without rounding, that for \(U\le V\),
\begin{equation}
 B(V)-B(U)=\frac{R(U,V)}{D(V)}.
 \label{eq:tail-code}
\end{equation}

The arithmetical theorem to be proved is the universal closure of
\begin{equation}
 \boxed{\quad
 N(q)\le U\ \wedge\ U\le V
 \quad\Longrightarrow\quad
 2^qR(U,V)\le D(V).
 \quad}
 \label{eq:ipra-target}
\end{equation}
All symbols in its matrix denote primitive recursive functions or
relations with primitive recursive characteristic functions. In particular,
its matrix is a quantifier-free formula after replacing order and bounded
tests by those characteristic functions. Formula~\eqref{eq:ipra-target}
therefore does not mean an informal quantification over infinite sums.

The order on natural numbers is decidable. Applying
\eqref{eq:ipra-target} to the smaller and larger of \(M,L\) gives
\eqref{eq:main-cauchy}. Conversely, the finite rational tail estimate proves
\eqref{eq:ipra-target} by multiplying by the positive integer \(2^qD(V)\).

\subsection{What convergence means here}

The sequence
\[
 b(q)=B\!\left(2^{\,2^{q+13}}\right)
\]
is a specified primitive recursive rational name. For \(r\ge q\),
\(\lvert b(r)-b(q)\rvert\le2^{-q}\).
In a constructive development of Cauchy reals, this name represents Brun's
constant. Its equivalence to the original sequence of partial sums follows
from the estimate for \emph{all} \(M,L\) beyond the cutoff, not merely for
the selected subsequence.

If the convention for a Cauchy modulus requires a strict inequality
\(<2^{-q}\), replace \(N(q)\) by \(N(q+1)\).
The proof never needs to decide whether there are finitely many or infinitely
many twin-prime pairs. A convergent series of nonnegative terms is compatible
with either possibility.

\subsection{Proof architecture}

The classical Selberg sieve provides the organizing idea
\cite{kedlaya,tau}, but we will prove the particular finite identity needed.
The steps are:
\begin{enumerate}
\item Count the forbidden residues of \(n(n+2)\) modulo squarefree integers.
\item Construct explicit rational coefficients \(\lambda_d\) whose squared
divisor sum majorizes the indicator of surviving the sieve.
\item Evaluate its main quadratic form as \(1/G(z)\), and bound the total
counting error by \(z^4\).
\item Prove \(G(s^2)\ge H_s^2/3\) by a finite prime-factor grouping.
\item Choose dyadic parameters and telescope a finite rational majorant
for the reciprocal tail.
\end{enumerate}
Sections~\ref{sec:finite}--\ref{sec:convergence} carry out these steps.
Section~\ref{sec:ipra} checks their induction and coding requirements.

\section{Elementary finite arithmetic}
\label{sec:finite}

We work with ordinary finite sets, sums and products as readable notation
for bounded loops. Empty sums are \(0\), empty products are \(1\), and
squarefreeness always refers to a \emph{positive} integer.
All rational denominators used below are proved positive.

\subsection{Division, prime factors, and finite coding}

For \(d>0\), division with remainder is obtained by a loop on the dividend.
Starting from quotient and remainder \((0,0)\), increase the remainder by
one at each step; when it reaches \(d\), reset it to \(0\) and increase
the quotient. Induction proves
\[
 n=dq+r,\qquad 0\le r<d.
\]
The stored pair can be coded by a primitive recursive pairing function.
No unbounded search is involved.

For \(n>1\), search the finite list \(2,\ldots,n\) for its first divisor.
The search succeeds because \(n\) is on the list. Its first result \(p\)
is prime: a divisor \(a\) of \(p\) with \(1<a<p\) would also divide \(n\)
and would have been found first. Divide out \(p\) and repeat, leaving the
state unchanged once the remaining factor is \(1\). Every successful step
strictly reduces that positive remaining factor, so a loop of \(n\) steps
suffices. Its invariant is
\[
 n=(\hbox{current remaining factor})
       \prod(\hbox{prime factors already recorded}).
\]
This constructs a prime-factor list, including multiplicities, by primitive
recursion.

For completeness, the divisibility fact used to prove uniqueness is
constructive. The subtraction version of the Euclidean algorithm maintains
the common divisors of a pair of nonnegative integers; whenever both
coordinates are positive, subtract the smaller from the larger. Their sum
strictly decreases, so a fixed loop of \(a+b\) steps, frozen at termination,
computes \(\gcd(a,b)\). Maintaining integer linear-combination coefficients
at the same time gives a Bezout identity. If \(p\) is prime and \(p\nmid a\),
then \(\gcd(p,a)=1\), so some explicitly computed integers \(u,v\) satisfy
\(up+va=1\). Multiplying by \(b\) shows
\[
 p\mid ab\quad\Longrightarrow\quad p\mid a\ \hbox{or}\ p\mid b.
\]
Here the case split \(p\mid a\) is a decidable divisibility test. Repeatedly
remove a matching prime from two factor lists to obtain uniqueness.
All the lists and coefficient traces have fixed finite bounds; a more
explicit coding audit is given in Section~\ref{sec:ipra}.

We may consequently use:
\[
 \rad(n)=\prod_{p\mid n}p,\qquad
 \omega(n)=\#\{p:p\mid n\}
 \quad(n\ge1),
\]
with \(\rad(1)=1\), and the Mobius function
\[
 \mu(n)=
 \begin{cases}
 (-1)^{\omega(n)},&n\text{ squarefree},\\
 0,&n\text{ not squarefree}.
 \end{cases}
\]
These are all obtained from bounded factorization. In particular,
\(\rad(n)\mid n\), and \(1\le\rad(n)\le n\).

\subsection{The residue counts}

For squarefree \(d\), define
\[
 \rho(d)=\#\{a:0\le a<d,\ d\mid a(a+2)\}.
\]
For a prime \(p\), the preceding divisibility fact gives the two possible
root residues \(0\) and \(-2\). They coincide exactly when \(p=2\).
Thus
\begin{equation}
 \rho(2)=1,\qquad
 \rho(p)=2\quad(p\text{ odd prime}),\qquad
 \rho(1)=1.
 \label{eq:local-rho}
\end{equation}

We recall the finite Chinese remainder bijection rather than treating it
as an existence principle. If \(a,b>0\) are coprime, compute an inverse
\(u\) of \(a\) modulo \(b\) from a Bezout identity. Given residues
\(0\le r<a\), \(0\le s<b\), set
\[
 x=r+a\bigl((s-r)u\bmod b\bigr).
\]
The remainder is read as its representative in \(\{0,\ldots,b-1\}\);
signed integers can be coded by pairs of natural numbers.
Then \(0\le x<ab\), \(x\equiv r\pmod a\), and \(x\equiv s\pmod b\).
If two representatives have these properties, their difference is
divisible by \(a\) and \(b\), hence by \(ab\): if the difference is \(ak\),
then \(b\mid ak\), and a Bezout identity for \(a,b\) implies \(b\mid k\).
Its absolute value is below \(ab\), so it is zero. This proves the bijection.

Applying the bijection successively to the distinct prime factors of \(d\)
gives
\begin{equation}
 \rho(d)=\prod_{p\mid d}\rho(p).
 \label{eq:rho-product}
\end{equation}
For squarefree \(d,e\), the local factors are at least \(1\), so
\begin{equation}
 \rho(\lcm(d,e))\le\rho(d)\rho(e).
 \label{eq:rho-lcm}
\end{equation}

\begin{lemma}[Finite counting error]
\label{lem:count-error}
For \(X\ge0\) and squarefree \(d\), let
\[
 A_d(X)=\#\{n:1\le n\le X,\ d\mid n(n+2)\}.
\]
Then
\begin{equation}
 \left|A_d(X)-X\frac{\rho(d)}d\right|\le\rho(d).
 \label{eq:count-error}
\end{equation}
\end{lemma}

\begin{proof}
Write \(X=qd+r\), \(0\le r<d\). The interval \(1,\ldots,X\)
consists of \(q\) complete consecutive blocks of \(d\) integers and a
remaining block of \(r\) integers. Each complete block contains exactly
one representative of each residue. A fixed root residue therefore occurs
either \(q\) or \(q+1\) times. Its difference from \(X/d=q+r/d\)
has absolute value at most \(1\). Sum this estimate over the
\(\rho(d)\) root residues and use the finite triangle inequality.
For \(d=1\) the error is exactly \(0\).
\end{proof}

\subsection{The cancellation identity}

\begin{lemma}[Squarefree Mobius cancellation]
\label{lem:mobius}
For squarefree \(v\),
\[
 \sum_{a\mid v}\mu(a)=
 \begin{cases}1,&v=1,\\0,&v>1.\end{cases}
\]
If \(r\mid v\) and \(v\) is squarefree, then
\begin{equation}
 \sum_{r\mid d\mid v}\mu(d)=
 \begin{cases}\mu(r),&v=r,\\0,&v\ne r.\end{cases}
 \label{eq:interval-mobius}
\end{equation}
\end{lemma}

\begin{proof}
Divisors of \(v\) correspond to subsets of its finite list of prime factors.
Expanding \(\prod_{p\mid v}(1-1)\) gives the first sum. For the second,
write \(d=ra\). Since \(v\) is squarefree, \(r\) and \(v/r\) are coprime,
so \(\mu(ra)=\mu(r)\mu(a)\) for \(a\mid v/r\). Apply the first identity.
The product expansion can also be proved directly by induction on the
length of the prime-factor list.
\end{proof}

\section{A finite quadratic sieve with explicit coefficients}
\label{sec:sieve}

Fix \(z\ge1\), and put
\[
 \sfD_z=\{d:1\le d\le z,\ d\text{ squarefree}\}.
\]
This set contains \(1\) and is closed under taking positive divisors.
For any squarefree \(d\), define the positive rationals
\begin{equation}
 g(d)=\frac{\rho(d)}d,\qquad
 h(d)=\prod_{p\mid d}\frac{\rho(p)}{p-\rho(p)}.
 \label{eq:gh}
\end{equation}
The prime denominators are positive: \(2-\rho(2)=1\), and
\(p-\rho(p)=p-2\ge1\) for odd primes.
By the empty-product convention \(g(1)=h(1)=1\).
Set
\[
 G=G(z)=\sum_{r\in\sfD_z}h(r)\ge1,\qquad
 W_d=\sum_{\substack{r\in\sfD_z\\d\mid r}}h(r).
\]
For \(d\in\sfD_z\), choose
\begin{equation}
 \boxed{\qquad
 \lambda_d=\frac{\mu(d)W_d}{g(d)G}.
 \qquad}
 \label{eq:weights}
\end{equation}
These are explicit finite rational expressions. There is no minimization
over real vectors in their definition.

\begin{lemma}[Normalization and coefficient bound]
\label{lem:weights}
The coefficients satisfy
\begin{equation}
 \lambda_1=1,\qquad
 |\lambda_d|g(d)=\frac{W_d}G\le1,\qquad
 |\lambda_d|\rho(d)\le d.
 \label{eq:coefficient-bound}
\end{equation}
For every \(r\in\sfD_z\),
\begin{equation}
 \sum_{\substack{d\in\sfD_z\\r\mid d}}\lambda_dg(d)
   =\frac{\mu(r)h(r)}G.
 \label{eq:transform}
\end{equation}
\end{lemma}

\begin{proof}
The normalization follows from \(W_1=G\).
Since \(d\) is squarefree, \(|\mu(d)|=1\), and \(0<W_d\le G\);
this proves the bounds. For the transform, expand \(W_d\) and change the
order of two finite sums:
\begin{align*}
 \sum_{\substack{d\in\sfD_z\\r\mid d}}\lambda_dg(d)
 &=\frac1G
   \sum_{\substack{v\in\sfD_z\\r\mid v}}
          h(v)\sum_{r\mid d\mid v}\mu(d)\\
 &=\frac{\mu(r)h(r)}G.
\end{align*}
The last equality is Lemma~\ref{lem:mobius}. All divisors occurring in the
inner sum belong to \(\sfD_z\), because \(v\le z\).
\end{proof}

\begin{lemma}[Quadratic identity]
\label{lem:quadratic}
The coefficients in \eqref{eq:weights} satisfy
\begin{equation}
 Q:=\sum_{d,e\in\sfD_z}\lambda_d\lambda_e
             g(\lcm(d,e))=\frac1{G(z)}.
 \label{eq:quadratic}
\end{equation}
\end{lemma}

\begin{proof}
For a prime \(p\), the definition of \(h\) gives
\[
 1+\frac1{h(p)}=\frac1{g(p)}.
\]
Expanding a finite product over the primes dividing a squarefree integer
\(v\) yields
\[
 \frac1{g(v)}
   =\prod_{p\mid v}\left(1+\frac1{h(p)}\right)
   =\sum_{r\mid v}\frac1{h(r)}.
\]
Multiplicativity on squarefree prime supports also gives
\[
 g(\lcm(d,e))=\frac{g(d)g(e)}{g(\gcd(d,e))}
   =g(d)g(e)\sum_{r\mid\gcd(d,e)}\frac1{h(r)}.
\]
Insert this into \(Q\) and interchange finite sums:
\begin{align*}
 Q
 &=\sum_{r\in\sfD_z}\frac1{h(r)}
       \left(\sum_{\substack{d\in\sfD_z\\r\mid d}}
                     \lambda_dg(d)\right)^2\\
 &=\sum_{r\in\sfD_z}\frac1{h(r)}
                 \left(\frac{\mu(r)h(r)}G\right)^2
   =\frac1{G^2}\sum_{r\in\sfD_z}h(r)
   =\frac1G.
\end{align*}
Here \(\mu(r)^2=1\) on \(\sfD_z\), and the transform is
\eqref{eq:transform}. This is an equality of finite rational expressions.
\end{proof}

\subsection{The sieve majorant and its error}

Let \(S(X,z)\) count integers \(1\le n\le X\) for which no prime at most
\(z\) divides \(n(n+2)\). This is a bounded condition.
For such an \(n\), the only member \(d\in\sfD_z\) dividing \(n(n+2)\)
is \(1\): any \(d>1\) has a prime divisor at most \(d\le z\).
Thus the following pointwise inequality holds for every \(n\):
\begin{equation}
 \mathbf 1_{\{n\text{ survives}\}}
 \le
 \left(\sum_{\substack{d\in\sfD_z\\d\mid n(n+2)}}\lambda_d\right)^2.
 \label{eq:majorant}
\end{equation}
If \(n\) survives, the right side is \(\lambda_1^2=1\); otherwise the
left side is \(0\) and a rational square is nonnegative.

Sum over \(1\le n\le X\), expand the square, and use the equivalence
\[
 d\mid n(n+2)\ \wedge\ e\mid n(n+2)
 \quad\Longleftrightarrow\quad
 \lcm(d,e)\mid n(n+2).
\]
We obtain
\[
 S(X,z)\le
 \sum_{d,e\in\sfD_z}\lambda_d\lambda_e A_{\lcm(d,e)}(X).
\]
By Lemma~\ref{lem:count-error}, write
\[
 A_v(X)=Xg(v)+E_v,\qquad |E_v|\le\rho(v)
\]
for the squarefree numbers \(v=\lcm(d,e)\).
There is no restriction \(v\le z\) here; indeed \(v\le z^2\).
Lemma~\ref{lem:quadratic} gives the main term \(X/G(z)\).
For the error, \eqref{eq:rho-lcm} and
\eqref{eq:coefficient-bound} give
\begin{align}
 \left|\sum_{d,e\in\sfD_z}\lambda_d\lambda_e E_{\lcm(d,e)}\right|
 &\le
 \sum_{d,e\in\sfD_z}|\lambda_d||\lambda_e|
                   \rho(d)\rho(e)\notag\\
 &=\left(\sum_{d\in\sfD_z}|\lambda_d|\rho(d)\right)^2
 \le\left(\sum_{d=1}^z d\right)^2
 \le z^4.
 \label{eq:sieve-error}
\end{align}
The final inequality uses \(z\) terms, each at most \(z\).
Consequently,
\begin{equation}
 S(X,z)\le\frac{X}{G(z)}+z^4.
 \label{eq:sifted}
\end{equation}

\begin{proposition}[Finite twin-prime bound]
\label{prop:finite-sieve}
For every \(X\ge0\) and \(z\ge1\),
\begin{equation}
 \boxed{\qquad
 T(X)\le z+\frac{X}{G(z)}+z^4.
 \qquad}
 \label{eq:finite-sieve}
\end{equation}
\end{proposition}

\begin{proof}
At most \(z\) pairs counted by \(T(X)\) have their smaller prime at most
\(z\). Every remaining pair has both primes greater than \(z\).
No prime at most \(z\) can divide either of those primes, and therefore
none divides their product. Its smaller member is counted by \(S(X,z)\).
Thus \(T(X)\le z+S(X,z)\), and \eqref{eq:sifted} proves the result.
\end{proof}

\section{The denominator estimate from a finite harmonic square}
\label{sec:denominator}

It remains to lower-bound \(G(z)\). This is the part where a standard
analytic presentation may invoke asymptotics for multiplicative sums.
The following direct inequality is sufficient here.

\begin{theorem}[Prime-support comparison]
\label{thm:harmonic}
For every \(s\ge1\), let \(H_s=\sum_{a=1}^s1/a\). Then
\begin{equation}
 \boxed{\qquad H_s^2\le3G(s^2).\qquad}
 \label{eq:harmonic-G}
\end{equation}
\end{theorem}

\begin{proof}
Expand the finite square:
\[
 H_s^2=\sum_{1\le a,b\le s}\frac1{ab}.
\]
Each ordered pair has the uniquely determined prime support
\(d=\rad(ab)\), which is squarefree and at most \(ab\le s^2\).
For \(d\in\sfD_{s^2}\), let
\[
 F_d(s)=
 \sum_{\substack{1\le a,b\le s\\\rad(ab)=d}}\frac1{ab}.
\]
We will show \(F_d(s)\le3h(d)\).

Write \(a=\prod_{p\mid d}p^{e_p}\),
\(b=\prod_{p\mid d}p^{f_p}\).
The exponents are nonnegative, and for each \(p\mid d\) they are not both
zero. Every exponent is at most \(s\): if \(p^e\mid a\le s\) with
\(e>0\), then \(e\le p^e\le s\); likewise for \(b\).
Uniqueness of factorization makes the map to the collection of exponent
pairs injective.

Enlarge the permitted exponent collections to the entire finite box
\[
 0\le e_p,f_p\le s,\qquad (e_p,f_p)\ne(0,0)
 \quad(p\mid d).
\]
This forgets the constraints \(a,b\le s\) and can only increase the sum
of nonnegative weights. Finite distributivity then gives
\begin{equation}
 F_d(s)\le
 \prod_{p\mid d}
 \left(
 \sum_{\substack{0\le e,f\le s\\(e,f)\ne(0,0)}}p^{-(e+f)}
 \right)
 =\prod_{p\mid d}\bigl(2A_p+A_p^2\bigr),
 \label{eq:box}
\end{equation}
where \(A_p=\sum_{e=1}^s p^{-e}\).
No infinite geometric series occurs in this step.

For the finite geometric sum, cancellation of consecutive terms gives
\[
 (p-1)A_p=1-p^{-s}\le1,\qquad
 0\le A_p\le\frac1{p-1}.
\]
If \(p\) is odd, hence \(p\ge3\), it follows that
\[
 2A_p+A_p^2
 \le\frac{2}{p-1}+\frac1{(p-1)^2}
 =\frac{2p-1}{(p-1)^2}
 \le\frac2{p-2}=h(p).
\]
The last inequality has positive denominators and reduces exactly to
\[
 2(p-1)^2-(2p-1)(p-2)=p\ge0.
\]
At \(p=2\), we instead have
\[
 2A_2+A_2^2\le3=3h(2).
\]
There is only one prime \(2\) in a squarefree support. Thus the product in
\eqref{eq:box} is at most \(h(d)\) when \(d\) is odd and at most
\(3h(d)\) when \(d\) is even. For \(d=1\), the product is empty and equals
\(1\), agreeing with the sole pair \(a=b=1\).
Hence in every case \(F_d(s)\le3h(d)\).

Finally partition the original finite pair set by \(d=\rad(ab)\):
\[
 H_s^2=\sum_{d\in\sfD_{s^2}}F_d(s)
       \le3\sum_{d\in\sfD_{s^2}}h(d)=3G(s^2).
\]
Every group with no original pair simply contributes \(0\).
\end{proof}

\begin{lemma}[Dyadic harmonic bound]
\label{lem:harmonic-dyadic}
For every \(t\ge0\),
\[
 H_{2^t}\ge1+\frac t2,\qquad
 G(2^{2t})\ge\frac{(t+2)^2}{12}.
\]
\end{lemma}

\begin{proof}
Separate the term \(1\) and the \(t\) blocks
\[
 2^{j-1}<a\le2^j,\qquad 1\le j\le t.
\]
The \(j\)-th block has \(2^{j-1}\) terms, each at least \(2^{-j}\);
its sum is at least \(1/2\). This proves the harmonic inequality by a
finite sum. Applying Theorem~\ref{thm:harmonic} with \(s=2^t\)
gives
\[
 G(2^{2t})\ge\frac{H_{2^t}^2}{3}
 \ge\frac{(1+t/2)^2}{3}
 =\frac{(t+2)^2}{12}.
\]
For \(t=0\), the calculation uses \(H_1=G(1)=1\), and remains valid.
\end{proof}

\section{The dyadic count and the explicit Cauchy modulus}
\label{sec:convergence}

\subsection{A uniform numerical bound}

\begin{lemma}
\label{lem:polynomial}
For every \(t\ge0\), \((t+1)^2\le2^{2t}\).
\end{lemma}

\begin{proof}
At \(t=0\) both sides are \(1\).
Since \(t+2\le2(t+1)\), squaring and using the induction hypothesis gives
\[
 (t+2)^2\le4(t+1)^2\le4\cdot2^{2t}=2^{2(t+1)}.
\]
This is induction on a quantifier-free numerical inequality.
\end{proof}

\begin{proposition}[Explicit dyadic counting theorem]
\label{prop:dyadic-count}
For every \(k\ge0\),
\[
 k^2T(2^k)\le1400\,2^k\le2048\,2^k.
\]
\end{proposition}

\begin{proof}
Put
\[
 X=2^k,\qquad t=k\mathbin{\mathrm{div}}10,\qquad z=2^{2t}.
\]
Division with remainder gives
\[
 10t\le k<10(t+1),\qquad
 k^2\le100(t+1)^2.
\]
Lemma~\ref{lem:harmonic-dyadic} shows that the main term in
Proposition~\ref{prop:finite-sieve} satisfies
\[
 \frac{k^2X}{G(z)}
 \le\frac{12k^2X}{(t+2)^2}
 \le1200\,\frac{(t+1)^2}{(t+2)^2}\,X
 \le1200X.
\]
For the error, Lemma~\ref{lem:polynomial} gives
\[
 k^2z^4\le100(t+1)^2\,2^{8t}
          \le100\,2^{10t}
          \le100\,2^k=100X.
\]
Since \(z\ge1\), also \(z\le z^4\), and so \(k^2z\le100X\).
Multiply \eqref{eq:finite-sieve} by \(k^2\) and add the three bounds:
\[
 k^2T(2^k)
 \le k^2z+\frac{k^2X}{G(z)}+k^2z^4
 \le(100+1200+100)X=1400X.
\]
Finally \(1400\le2048=2^{11}\).
This proof includes \(t=0\), and even \(k=0\); it never divides by \(k\).
For \(k\ge1\), division by \(k^2\) yields the familiar form
\[
 T(2^k)\le2048\,\frac{2^k}{k^2}.
\]
\end{proof}

\subsection{A finite tail estimate}

Write \(C=2048\). For a dyadic shell indexed by \(j\ge1\), put
\[
 V_j=\sum_{\substack{2^j\le p<2^{j+1}\\\Twin(p)}}
              \left(\frac1p+\frac1{p+2}\right).
\]
Each summand is at most \(2/p\le2/2^j\).
The number of summands is at most \(T(2^{j+1})\). The inclusive endpoint
in \(T\) only enlarges this bound. Proposition~\ref{prop:dyadic-count}
therefore gives
\begin{equation}
 V_j\le\frac2{2^j}T(2^{j+1})
       \le\frac{4C}{(j+1)^2}.
 \label{eq:shell}
\end{equation}

For \(j\ge1\), all denominators below are positive, and
\begin{equation}
 \frac1{(j+1)^2}\le\frac1{j(j+1)}
  =\frac1j-\frac1{j+1}.
 \label{eq:telescope-step}
\end{equation}
Summing only from \(j=m\) to a finite \(K\ge m\) gives
\begin{equation}
 \sum_{j=m}^{K}\frac1{(j+1)^2}
 \le\frac1m-\frac1{K+1}\le\frac1m.
 \label{eq:telescope}
\end{equation}
This uses no assertion about the convergence of an auxiliary infinite
series.

\begin{proposition}[Finite reciprocal tail]
\label{prop:tail}
For \(m\ge1\) and \(L\ge0\),
\[
 \sum_{\substack{2^m\le p<L\\\Twin(p)}}
          \left(\frac1p+\frac1{p+2}\right)
 \le\frac{4C}{m}=\frac{8192}{m}.
\]
\end{proposition}

\begin{proof}
If \(L\le2^m\), the sum is empty. Otherwise choose, for instance,
\(K=m+L\). The elementary inequality \(L<2^{K+1}\) follows from
\(u<2^{u+1}\), proved by quantifier-free induction.
Every integer \(p\) with \(2^m\le p<L\) lies in a unique shell
\([2^j,2^{j+1})\) with \(m\le j\le K\).
To find its index, use the bounded search for the first
\(j\in\{m,\ldots,K\}\) with \(p<2^{j+1}\); the previous endpoint,
or \(p\ge2^m\) in the first case, gives \(2^j\le p\).
Thus the original sum is at most \(\sum_{j=m}^K V_j\).
Equations~\eqref{eq:shell} and \eqref{eq:telescope} give the claimed bound.
\end{proof}

\subsection{Convergence with a primitive recursive modulus}

Fix \(q\ge0\), and set
\[
 m=2^{q+13},\qquad N(q)=2^m.
\]
Then \(m\ge1\) and
\begin{equation}
 \frac{8192}{m}
 =\frac{2^{13}}{2^{q+13}}
 =2^{-q}.
 \label{eq:modulus-arithmetic}
\end{equation}
Let \(M,L\ge N(q)\). By decidable comparison, assume first that \(L\le M\).
Nonnegativity of the summands and Proposition~\ref{prop:tail} give
\[
 0\le B(M)-B(L)
 =\sum_{\substack{L\le p<M\\\Twin(p)}}
       \left(\frac1p+\frac1{p+2}\right)
 \le\sum_{\substack{2^m\le p<M\\\Twin(p)}}
       \left(\frac1p+\frac1{p+2}\right)
 \le2^{-q}.
\]
If \(M<L\), apply the same calculation with the names interchanged.
This proves Theorem~\ref{thm:main}(iii).

For the precise integer version, take \(N(q)\le U\le V\), use
\eqref{eq:tail-code}, and multiply the tail inequality by
\(2^qD(V)>0\). The result is exactly \eqref{eq:ipra-target}.
This completes the finite mathematical argument.

\begin{remark}[Effectivity and numerical usefulness]
The formula \(N(q)=2^{2^{q+13}}\) is a composition of primitive recursive
exponentiation functions. Composing it with bounded primality testing,
factorial, and finite summation remains primitive recursive.
Its size is immaterial to primitive recursiveness, although it makes this
particular approximation algorithm impractical.
The theorem supplies a certified convergence rate, not a competitive
method for computing many digits of Brun's constant.
\end{remark}

\section{Why the proof belongs to iPRA}
\label{sec:ipra}

The mathematical argument above uses only rational arithmetic and finite
combinatorics. That observation alone would not establish a weak-theory
result: a true statement about a terminating computation can require a
stronger theory to prove uniformly. We now specify the definitions and
induction invariants that justify the iPRA claim.

\subsection{The underlying theory}

Use ordinary first-order intuitionistic logic with equality and one
function symbol for each primitive recursive definition. The arithmetic
axioms are the defining equations for zero, successor, projections,
composition and primitive recursion, successor injectivity,
\(Sx\ne0\), and induction restricted to quantifier-free formulas.
Parameters are allowed in the induction schema.

This is exactly the presentation defined by \code{PRAAx} in the inspected
ConciseIPRA source \cite{proveit}; \code{PRA.qfInd} exposes its induction
rule. Thus we do not silently replace the requested theory by Heyting
arithmetic with induction for all formulas.

\subsection{Arithmetic and rational codes}

Addition and multiplication are defined by primitive recursion.
Truncated subtraction, comparison, parity, bounded quotient and remainder,
finite sums, products, factorial, and exponentiation are primitive recursive.
The usual ordered semiring identities follow by quantifier-free induction
on their defining recursions. For instance, distributivity is proved by
induction on the recursion argument of multiplication, with the other
arguments as free parameters.

An integer can be represented by a pair \((a,b)\) denoting \(a-b\).
Equality is \(a+d=c+b\) for \((a,b)\) and \((c,d)\).
Addition and multiplication have polynomial definitions:
\[
 (a,b)+(c,d)=(a+c,b+d),\qquad
 (a,b)(c,d)=(ac+bd,ad+bc).
\]
Negation swaps the entries. Order is expressed by a natural-number
comparison. This permits all signed calculations involving \(\mu\),
\(\lambda_d\), and the errors \(E_v\) without adjoining negative numbers
as a new domain.

A rational is an integer numerator code and a positive natural denominator.
Reduction to lowest terms is unnecessary. For denominators \(b,d>0\),
\[
 \frac ab+\frac cd=\frac{ad+bc}{bd},\qquad
 \frac ab\frac cd=\frac{ac}{bd},\qquad
 \frac ab\le\frac cd\ \Longleftrightarrow\ ad\le cb.
\]
Inverting a nonzero rational swaps the absolute numerator and denominator
and retains its sign. The proof uses inverses only with a separately proved
positive input. In particular \(g(d)>0\), \(h(d)>0\), \(G(z)\ge1\);
the factors \(p-\rho(p)\) are positive.
All denominator clearing is therefore legitimate in the integer coding.
Rational squares are nonnegative by a decidable sign comparison followed
by the integer multiplication identities.

Finite rational sums are recursive folds of these explicit operations.
Their algebraic rules are proved by induction on fold length.
For example, double-sum interchange is obtained by induction on one
length and then on the other, using distributivity and associativity at
each step. Nonnegative-term comparison is a fold induction on the
partial-sum inequality. These induction formulas compare particular
integer codes, and are quantifier-free.

\subsection{Bounded conditions can be put into induction invariants}

For a primitive recursive Boolean predicate \(P(i,\vec a)\), define
\[
 \operatorname{All}_P(0,\vec a)=1,\qquad
 \operatorname{All}_P(n+1,\vec a)
   =\operatorname{All}_P(n,\vec a)\cdot P(n,\vec a).
\]
Its value is \(1\) exactly when all entries \(i<n\) pass the test.
The analogous bounded existential test is obtained by testing whether the
sum of the Boolean entries is positive. Bounded searches and bounded
maxima are defined by folds with a success flag. For an existential test,
the fold also records the first successful index. Its invariant says that
if the flag is \(1\), the stored index is below the processed bound and
passes \(P\). This supplies the witness constructively.

The facts required to use these encodings are themselves proved by
quantifier-free induction. A useful explicit strengthened invariant for
recovering entries is
\[
 \operatorname{All}_P(n,\vec a)=1\ \wedge\ i<n
       \ \Longrightarrow\ P(i,\vec a)=1,
\]
with \(i,\vec a\) as free parameters; the induction step splits the
decidable cases \(i<n\), \(i=n\), and \(i>n\).
In the other direction, under assumptions supplying the bounded entries,
the defining multiplication recurrence proves that the aggregate is \(1\):
with the final bound \(b\) fixed, induct on the quantifier-free assertion
\(j\le b\Rightarrow\operatorname{All}_P(j,\vec a)=1\).
This establishes the familiar bounded-quantifier notation without
induction on an unbounded universal formula.

For a nested bounded formula, apply this construction from the innermost
test outward. Its characteristic function remains primitive recursive,
and its use inside an induction can be replaced by equality of that
function to \(1\). The finite assertions in the present proof involve
only such bounded conditions and coded arithmetic equalities or
inequalities.

\subsection{Explicit bounds for the finite objects}

Here are bounds sufficient to make the construction literal.
They need not be efficient.
\begin{itemize}
\item The least-divisor search on \(n>1\) tests at most \(n\) candidates.
Factorization at \(0\) is assigned a fixed default, and its correctness
claims concern positive inputs. The factorization loop has \(n\) stages.
A padded factor trace has length
\(n\), and entries at most \(n\); it can be stored as digits in base \(n+1\),
together with the length \(\ell\le n\) of its active prefix. Explicitly,
the code is \(\operatorname{digits}+\ell(n+1)^n\), with
\(\operatorname{digits}<(n+1)^n\), hence is below \((n+1)^{n+1}\).
\item The subtraction Euclidean algorithm has at most \(a+b\) active
stages. Its Bezout coefficients start with absolute value at most \(1\)
and at each stage are copied or subtracted. Their absolute values are
bounded by \(2^{a+b+1}\). Sign and magnitude codes give a bounded trace.
\item Residue tables modulo \(v\) have \(v\) entries. In the sieve,
\(d,e\le z\) and \(\lcm(d,e)\le z^2\), so all such tables have at most
\(z^2\) entries. Their truth values are computed by remainder tests.
\item Divisor and squarefree tests at \(d\le z\) use ranges bounded by
\(z\). The sums defining \(G,W_d,\lambda_d\) have at most \(z\) terms.
The quadratic form has at most \(z^2\) terms.
\item In the harmonic comparison, \(ab\le s^2\).
For each \(d\le s^2\), there are at most \(s^2\) prime slots and
two exponents in \(\{0,\ldots,s\}\) at each slot.
Thus \((s+1)^{2s^2}\) bounds the number of padded exponent codes.
Their products and weights are computed by bounded folds.
\item A dyadic shell index for the tail below \(L\) is found in the
explicit range \(m,\ldots,m+L\). No logarithm function or unbounded
least-witness principle is required.
\end{itemize}
The coding of an entire rational computation is likewise a primitive
recursive trace: each next numerator and denominator is computed by the
displayed operations. If a bound on trace entries is wanted in an
invariant, compute the maximum along that same finite trace by a fold.
There is no existentially postulated sequence of intermediate rationals.

\subsection{The induction audit}

Every induction used in the proof can be organized as one of the following
finite invariant checks. Ambient bounds such as \(n,z,s,X,m,L\) are fixed
parameters during the relevant induction.

\renewcommand{\arraystretch}{1.12}
\begin{longtable}{@{}p{0.21\textwidth}p{0.28\textwidth}p{0.44\textwidth}@{}}
\toprule
Part of proof & Induction variable & Quantifier-free invariant after coding\\
\midrule
\endhead
Division & Number of processed increments &
Dividend \(=\) divisor \(\times\) quotient \(+\) remainder, with remainder
below divisor.\\
Factorization & Bounded stage counter &
Original number \(=\) remaining factor \(\times\) recorded product;
the recorded factors pass a bounded primality test; active steps decrease
the remaining factor.\\
Bezout and gcd & Bounded stage counter &
The two stored coordinates have the stored linear-combination
representations; common divisors up to \(a+b\) agree with the initial
pair; either the computation has stopped or its sum has decreased.\\
Factor uniqueness & Length of cancelled prefix &
Matching prime factors have been removed, and the two remaining
products agree; all residual factors are prime.\\
Chinese remainders & Length of prime prefix &
The computed reconstruction and residue projection are inverse on their
bounded domains; the product of root counts equals the root-table count.\\
Mobius cancellation & Length of prime prefix &
The partial subset sum equals the corresponding product of local
\((1-1)\) factors.\\
Finite algebra & Sum/product length &
The two accumulated rational codes are equal, or one accumulated code
is at most the other.\\
Selberg transform & Divisor-sum length &
The accumulated reindexed sums agree; the interval cancellation identity
holds for indices bounded by \(z\).\\
Quadratic identity & Bounded sum lengths &
The partial quadratic expansion equals its reindexed divisor expression;
the completed expansion equals \(1/G\).\\
Harmonic comparison & Prime-prefix and box counters &
Partial exponent-box sums factor; original pair codes map injectively to
box codes; all weights are nonnegative.\\
Dyadic estimates & \(t\) or shell counter &
\(H_{2^t}\ge1+t/2\), \((t+1)^2\le4^t\), or the displayed finite
telescoping inequality.\\
\bottomrule
\end{longtable}

For the termination statements in this table, a numerical strengthening
removes any unbounded termination assertion. For the factor loop at stage
\(i\le n\), either it has reached \(1\), or its remaining factor plus \(i\)
is at most \(n\). At stage \(n\) the latter case contradicts positivity.
For the subtraction Euclidean loop, either it has stopped, or the
coordinate sum plus the stage counter is at most the initial \(a+b\).
At stage \(a+b\) two positive coordinates are impossible.
Both strengthened assertions are quantifier-free after the state is coded.

In the factor-uniqueness and CRT rows, statements about all entries are
bounded by the fixed input products or trace lengths. Their bounded
Boolean aggregates can be carried in the invariant. The elementary
preservation calculations are those already given in
Section~\ref{sec:finite}; no induction over all arbitrary formulas or all
unbounded witnesses is needed.

\subsection{The resulting derivability claim}

Expand the finite definitions using their primitive recursive symbols,
interpret integer and rational notation by the preceding codes, and
replace bounded conditions in induction invariants by their characteristic
functions. The proofs in Sections~\ref{sec:finite}--\ref{sec:convergence}
then use the defining equations, intuitionistic reasoning with decidable
arithmetic cases, and the quantifier-free inductions just listed.
After the final denominator multiplication they derive
\eqref{eq:ipra-target}; universal introduction closes the free parameters
\(q,U,V\).

Only the finitely many specific primitive recursive definitions used by
this proof are involved. No primitive recursive universal interpreter for
all primitive recursive programs is assumed.

This is a human-readable proof of the iPRA theorem, together with its
translation discipline. It is not a claim that a formal derivation
certificate has already been generated. Producing that certificate is
the object-theory stage of the Lean plan below.

In particular, the proof requires neither an infinite Euler product,
a prime-distribution asymptotic, a compactness or completeness principle,
nor a choice of an unknown convergence threshold. All numerical case
splits concern decidable tests. There is no use of excluded middle for an
arbitrary proposition, Markov's principle, or a decision about infinitely
many twin primes.

\section{Repository reuse and Lean formalization}
\label{sec:lean}

\subsection{The inspected foundation}

The ProveIt repository was inspected at the fixed commit
\[
 \texttt{\small 8c135f6f2a9d041cde785a1a632c44ad98065b20}.
\]
The snapshot's toolchain and Mathlib dependency are both pinned to
version \(4.32.0\). The relevant source directory is
\begin{center}
\code{Logic/PrimitiveRecursiveArithmetic/ConciseIPRA/Lean/ConciseIPRA/}.
\end{center}
The following declarations were read in source, rather than inferred
from file names or an informal summary:
\begin{itemize}
\item \code{PR.lean}: the types \code{PR} and \code{PTm}, the axiom
family \code{PRAAx}, and the rule \code{PRA.qfInd}.
\item \code{PRAToScalar.lean}: \code{FromPRA.code},
\code{FromPRA.interp}, and proofs translating every arithmetic axiom,
including the quantifier-free induction axiom.
\item \code{BiInterpretation.lean}: \code{biInterp},
\code{pra\_iff\_teq}, \code{teq\_iff\_pra}, and the round-trip theorems.
\item \code{NativeBiInterpretation.lean}: \code{NFm.ofPRA},
\code{pra\_iff\_native}, and \code{native\_iff\_pra}.
\end{itemize}
These sources establish the existing syntactic translation interfaces
\cite{proveit}. They do not already provide the finite sieve proved here.
The foundation was not rebuilt locally in producing this package.

The inspected draft
\code{docs/drafts/prime\_reciprocal\_ipra\_2026-11-15.zip}
contains useful bounded-factorization and harmonic-sum methodology
\cite{prime-draft}. Its Lean scaffold contains explicit unproved
placeholders. It is therefore background for the proof design, not an
importable completed number-theory proof. The present argument proves
all of its required arithmetic facts independently.

\subsection{Two distinct formalization goals}

There are two layers to complete.

\paragraph{Finite mathematical layer.}
Prove the rational finite theorems in Lean, using computable finite lists
or finite sets. The primary targets are the quadratic identity, the
harmonic comparison, the dyadic count, and the finite tail bound.
Lean's rational numbers are convenient for checking these identities.
No infinite sums or real analysis are needed.

\paragraph{Object-theory layer.}
Build actual primitive recursive syntax codes for the displayed numerical
functions and a derivation in the deeply embedded theory \code{PRAAx}.
The final target has the form
\[
 []\ \vdash_{\mathsf{PRAAx}}\ \mathrm{BrunFormula}.
\]
A Lean theorem asserting the ordinary truth of
\eqref{eq:ipra-target} is valuable, but it does not on its own establish
that iPRA proves it. The additional syntactic derivation is essential.

\subsection{Recommended modules and interfaces}

\renewcommand{\arraystretch}{1.12}
\begin{longtable}{@{}p{0.30\textwidth}p{0.63\textwidth}@{}}
\toprule
Proposed module & Completion target\\
\midrule
\endhead
\code{Brun/FiniteArithmetic} &
Prove correctness of the bounded prime test; construct the factor trace,
squarefree support, divisor list and rational codes.\\
\code{Brun/Residues} &
Prove the explicit CRT reconstruction, \(\rho(d)=\prod_{p\mid d}\rho(p)\),
and the rational error bound for \(A_d(X)\).\\
\code{Brun/SelbergIdentity} &
Define \(G,W,\lambda\); prove \(\lambda_1=1\),
\(|\lambda_d|\rho(d)\le d\), the transform, and \(Q=1/G\).
Then derive \(T(X)\le z+X/G+z^4\).\\
\code{Brun/DenominatorLower} &
Construct the injection from bounded pairs into prime exponent boxes;
prove finite product factorization and \(H_s^2\le3G(s^2)\).\\
\code{Brun/DyadicCount} &
Prove the harmonic and polynomial inequalities and instantiate
\(t=k\mathbin{\mathrm{div}}10\), obtaining \(k^2T(2^k)\le2048\,2^k\).\\
\code{Brun/Tail} &
Construct the bounded shell index, bound each shell, telescope the
finite sum, and prove \eqref{eq:main-tail}.\\
\code{Brun/Cauchy} &
Finish the supplied conditional assembly, and prove equivalence with
the factorial-denominator integer statement.\\
\code{Brun/IPRA/Codes} &
Construct explicit \code{PR} syntax trees and defining-equation
derivations for all operations used by the proof.\\
\code{Brun/IPRA/Derivation} &
Internalize the lemmas using \code{PRA.qfInd}, with a quantifier-free
syntactic witness at every arithmetic induction, and derive the closed
formula.\\
\code{Brun/IPRA/Native} &
Apply the existing faithful interpretations to the completed
\code{PRAAx} derivation.\\
\bottomrule
\end{longtable}

The most substantial new tasks are the finite prime-support injection and
the syntactic arithmetic library. The number-theoretic bottleneck is
specified by Theorem~\ref{thm:harmonic}; no unknown analytic estimate is
left for the formalization to discover.

\subsection{A concrete semantic specification}

The supplied \code{lean/BrunSpecification.lean} defines:
\begin{itemize}
\item \code{primeB} and \code{twinB} by bounded list scans;
\item \code{brunPartial}, with the strict endpoint \(p<N\);
\item \code{sieveConstant}, \code{exponentIndex}, and \code{cutoff},
with constants exactly as in this article;
\item \code{tailNumerator} and \code{ArithmeticBrunClaim},
which encode \eqref{eq:integer-tail} and \eqref{eq:ipra-target};
\item \code{DyadicTailClaim} and \code{BrunCauchyClaim};
\item candidate conditional assembly proofs showing how a proof of the
finite tail bound yields the stated Cauchy estimate.
\end{itemize}
There are no newly declared axioms or \code{sorry} placeholders in that
file. The missing sieve theorem is a named proposition supplied as an
\emph{explicit hypothesis} to the assembly lemma; it is not asserted as
proved. The source has not been compiled because a Lean/Lake executable
was unavailable in the authoring environment.

The crucial assembly principle is elementary. Suppose every \(a(n)\)
beyond a cutoff lies in the same interval
\([a(N),a(N)+\epsilon]\). Then any two such values differ in absolute value
by at most \(\epsilon\), not \(2\epsilon\). This explains why the tail
bound produces the stated modulus without an additional factor of two.

\subsection{Building the object-language formula}

The existing \code{PR} syntax is generated by zero, successor,
projections, composition and primitive recursion. Define within this syntax
the fixed operations needed for the proof: comparison, quotient,
remainder, exponentiation, factorial, bounded folds, twin testing,
the cutoff, and the integer tail numerator.

For each function, prove its computation equations in \code{PRAAx}.
For each finite lemma, prove the corresponding coded statement by the
specific invariant in Section~\ref{sec:ipra}. Algebraic normalization can
be automated by a proof-producing macro that emits uses of the semiring
and rational-code lemmas. A computation that merely reports a true
Boolean is not a substitute for a uniform derivation.

With variables ordered as \(q,U,V\), form a primitive recursive Boolean
\(b(q,U,V)\) that is \(1\) precisely when the antecedent of
\eqref{eq:ipra-target} is false or its conclusion is true.
Both tests are decidable. The equality \(b(q,U,V)=1\) is an atomic
formula \code{brunMatrix : Fm PTm 3}; its \code{Fm.QF} evidence is immediate.
Universally close the three variables to obtain
\code{brunFormula : Fm PTm 0}.

The final task is to construct a term of
\[
 \code{Prf}\ \mathsf{PRAAx}\ []\ \mathrm{brunFormula},
\]
using the repository's exact notation and parameter order.
All arithmetic induction applications must go through
\code{PRA.qfInd} or a derived macro that supplies the required
\code{Fm.QF} witness.
Structural induction in Lean on a syntax tree or on a derivation is
metatheory; it does not grant unrestricted numerical induction in the
object theory.

\subsection{Transfer and verification gates}

After a derivation \code{brunPRA} is constructed, its scalar translation
is obtained through \code{FromPRA.interp.sound}; the native translation is
available through \code{pra\_iff\_native}. Thus there is no additional
bi-interpretability theorem to prove for this number-theory application.

Completion should be checked in the pinned ProveIt environment:
\begin{enumerate}
\item Compile the finite arithmetic modules and the supplied specification.
\item Compile the \code{PR} codes and the actual
\code{PRAAx} derivation.
\item Compile the scalar and native corollaries.
\item Inspect \code{\#print axioms} for the final declarations and verify
that there is no \code{sorryAx} or assumption standing for Brun's theorem.
\item Check that the final result is a derivation in \code{PRAAx}, not
only a theorem about standard natural numbers.
\item Review the \code{Fm.QF} witness at every object-level induction.
\end{enumerate}
The accompanying \code{FORMALIZATION\_PLAN.md} gives the same route in a
form convenient for implementation. No successful Lean build is claimed
for this package.

\section{Artifacts, verification, and scope of the result}
\label{sec:artifacts}

\subsection{What is in the archive}

The package contains this article in LaTeX source and compiled PDF, the
Lean specification and implementation plan, and exact-arithmetic Python
diagnostics. It also records the inspected repository snapshot and
reproduction commands. The proof does not depend on an unrecorded
calculation or a numerical table.

\begin{center}
\begin{tabular}{@{}>{\raggedright\arraybackslash}p{0.40\textwidth}>{\raggedright\arraybackslash}p{0.52\textwidth}@{}}
\toprule
File or directory & Purpose\\
\midrule
\code{brun\_ipra.tex}, \code{sections/} &
Complete human-readable argument and the iPRA induction audit.\\
\code{brun\_ipra.pdf} &
Compiled article.\\
\code{lean/BrunSpecification.lean} &
Uncompiled specification and candidate conditional assembly proofs.\\
\code{lean/FORMALIZATION\_PLAN.md} &
Detailed semantic and object-theory implementation route.\\
\code{lean/SOURCE\_LEDGER.md} &
Exact source paths, commit and blob identifiers, toolchain and status.\\
\code{artifacts/finite\_sieve.py} &
Exact finite calculations of the coefficients, quadratic form,
prime-support grouping and partial sums.\\
\code{tests/test\_finite\_sieve.py} &
Ten unit tests checking the implemented finite identities.\\
\code{artifacts/verification.json} &
Reproducible diagnostic results: 1,730 exact arithmetic assertions.\\
\code{README.md}, \code{Makefile} &
Reading guide and reproduction commands.\\
\bottomrule
\end{tabular}
\end{center}

\subsection{What the diagnostics check}

The program uses Python integers and \code{fractions.Fraction}; it uses no
floating-point comparisons. The checks include:
\begin{itemize}
\item prime-factor reconstruction, Mobius cancellation, and root-residue
counts for squarefree moduli;
\item the counting error, \(\lambda_1=1\), the Mobius transform,
\(Q=1/G\), and the bound \(|\lambda_d|\rho(d)\le d\);
\item the partition of \(H_s^2\) by \(\rad(ab)\), and the individual
prime-support group bounds;
\item the square majorant, total sieve error, and twin-prime inclusion;
\item the harmonic and elementary exponential inequalities on sample
parameters;
\item agreement of exact finite Brun sums with factorial-denominator
integer codes, including the double occurrence of \(5\).
\end{itemize}

All ten unit tests passed. The separate JSON diagnostic records 1,730
assertions. These are finite consistency checks on the transcription and
implementation. The universal statements are proved by the argument in
this article; the finite tests do not certify those universal statements,
and do not constitute a Lean proof.

For example, at \(z=5\),
\[
 \sfD_5=\{1,2,3,5\},\qquad G(5)=\frac{14}{3},
\]
and the coefficients are
\[
 \lambda_1=1,\quad
 \lambda_2=-\frac37,\quad
 \lambda_3=-\frac9{14},\quad
 \lambda_5=-\frac5{14}.
\]
The program evaluates the quadratic form exactly as \(3/14=1/G(5)\).
These modest examples are useful checks of signs and normalization,
where a transcription error could invalidate the sieve.

\subsection{Reproduction}

From the extracted package root:
\begin{lstlisting}
python3 -m unittest discover -s tests -v
python3 artifacts/finite_sieve.py --json artifacts/verification.json
latexmk -pdf -interaction=nonstopmode -halt-on-error brun_ipra.tex
\end{lstlisting}
The Python code needs only the standard library.
The PDF build needs a usual LaTeX installation with the packages listed
in the preamble.
Lean compilation must be performed in a workspace providing the pinned
Lean/Mathlib environment, for example:
\begin{lstlisting}
lake env lean /absolute/path/to/lean/BrunSpecification.lean
\end{lstlisting}
That command is a reproduction instruction, not a record of a build
performed here.

\subsection{Final mathematical and formal status}

The completed result is the finite theorem
\eqref{eq:ipra-target}, with its full human-readable sieve proof,
an explicit primitive recursive Cauchy modulus, and a detailed audit of
the permitted induction. It expresses constructive convergence of the
original pair reciprocal series.

The implementation work still needed for a full machine-checked result is
precisely the finite sieve library and the object-language derivation
described in Section~\ref{sec:lean}. The archive does not present that
work as already completed.


\begin{thebibliography}{9}
\raggedright
\bibitem{brun}
V.~Brun, \emph{La s\'erie \(1/5+1/7+1/11+1/13+\cdots\), o\`u les
d\'enominateurs sont nombres premiers jumeaux est convergente ou finie},
Bulletin des Sciences Math\'ematiques 43 (1919), 100--104, 124--128.
Historical reference; the proof in this article is supplied independently.

\bibitem{kedlaya}
K.~S.~Kedlaya, \emph{Notes on analytic number theory}, Sections 13 and 14,
\emph{The Selberg sieve} and \emph{Applying the Selberg sieve}.
\url{https://kskedlaya.org/ant/chap-selberg.html};
\url{https://kskedlaya.org/ant/chap-selberg2.html}.
Accessed 15 November 2026.

\bibitem{tau}
Tel Aviv University sieve theory course, \emph{Selberg's sieve:
bounding the number of twin primes}, course notes dated 7 May 2015,
hosted by Z.~Rudnick.
\url{https://www.math.tau.ac.il/~rudnick/courses/sieves2015/selberg%20sieve%20twin%20primes.pdf}.
The finite prime-support argument here is a direct replacement for the
infinite geometric expansions appearing in standard presentations.

\bibitem{proveit}
V.~Reshetnikov and contributors, \emph{ProveIt: ConciseIPRA},
repository snapshot
\texttt{8c135f6f2a9d041cde785a1a632c44ad98065b20},
15 November 2026.
\url{https://github.com/VladimirReshetnikov/ProveIt/tree/8c135f6f2a9d041cde785a1a632c44ad98065b20/Logic/PrimitiveRecursiveArithmetic/ConciseIPRA}.
Exact source paths, blob identifiers and inspection status are recorded in
\code{lean/SOURCE\_LEDGER.md}.

\bibitem{prime-draft}
\emph{Primitive-recursive divergence of the prime reciprocal series},
research package \code{prime\_reciprocal\_ipra\_2026-11-15.zip},
in \code{ProveIt/docs/drafts}, at the same snapshot.
Used as methodological background for bounded factor traces and arithmetic
encodings, not as an already completed Lean theorem.
\end{thebibliography}
\end{document}
